Fix \(\mathbb P\in\mathcal D_{d,\epsilon}^{\mathrm{obs}}\), write \(q_{\jmath,x}=\mathbb P\{(A,Y)=\jmath,X=x\}\), put \(L=L_d\), \(m=n/8\), and set \(\tau=L/m\).  We work in the uncapped comparison experiment, namely with an infinite continuation of the independent observations allowed in \(\mathrm{ass:iid\mbox{-}sampling}\).  Since \(M\sim\operatorname{Pois}(n/4)=\operatorname{Pois}(2m)\), the joint probability generating function of all pilot and evaluation counts is
\[
 \begin{split}
 &E_{\mathbb P}\!\left[
   \prod_{x,\jmath}z_{\jmath,x}^{N'_{\jmath,x}}
                      w_{\jmath,x}^{N_{\jmath,x}}
 \right]\\
 &\quad=\exp\!\left[
 2m\left\{\sum_{x,\jmath}q_{\jmath,x}
       \frac{z_{\jmath,x}+w_{\jmath,x}}2-1\right\}
 \right]\\
 &\quad=\prod_{x,\jmath}
   \exp\{mq_{\jmath,x}(z_{\jmath,x}-1)\}
   \exp\{mq_{\jmath,x}(w_{\jmath,x}-1)\}.
 \end{split}                                                     \tag{1}
\]
Thus the displayed counts are mutually independent, \(N'_{\jmath,x}\sim\operatorname{Pois}(mq_{\jmath,x})\), \(N_{\jmath,x}\sim\operatorname{Pois}(mq_{\jmath,x})\), and the entire pilot array is independent of the evaluation array.  This derives the required Poisson splitting directly from the sampling assumption and the construction in \(\mathrm{def:jackson\mbox{-}factorial\mbox{-}estimator}\).

For a scalar \(N\sim\operatorname{Pois}(mq)\), the definition of \(U_h\) and the binomial theorem give the identity of formal power series
\[
 \sum_{h=0}^{\infty}U_h(N;z)\frac{w^h}{h!}
 =e^{-zw}(1+w/m)^N.                                                \tag{2}
\]
Because every coefficient in (2) is a finite polynomial in \(N\), coefficient extraction after expectation is legitimate.  The Poisson generating function yields
\[
 E[e^{-zw}(1+w/m)^N]=e^{(q-z)w}                                   \tag{3}
\]
and, with two formal variables \(w,v\),
\[
 E\!\left[e^{-z(w+v)}(1+w/m)^N(1+v/m)^N\right]
 =\exp\{(q-z)(w+v)+(q/m)wv\}.                                    \tag{4}
\]
The coefficient of \(w^h/h!\) in (3) is \((q-z)^h\).  Expanding the last exponential in (4), the coefficient of \(w^hv^t/(h!t!)\) is
\[
 \sum_{\ell=0}^{h\wedge t}
 \binom h\ell\binom t\ell\ell!(q/m)^\ell
 (q-z)^{h+t-2\ell}.                                               \tag{5}
\]
Equations (3)--(5), together with the pilot--evaluation independence in (1), prove both centered-factorial identities conditional on the pilot counts.

We next prove the local pilot bounds used by the Jackson branch of the estimator.  Fix \(x\), abbreviate \(q_\jmath=q_{\jmath,x}\), \(c_\jmath=N'_{\jmath,x}/m\), and let \(h_\jmath,\ell_\jmath,u_\jmath,b_\jmath,r_\jmath\) be the quantities in \(\mathrm{def:jackson\mbox{-}factorial\mbox{-}estimator}\).  Define
\[
 G_x:=\bigcap_{\jmath\in\mathcal J}
 \{|c_\jmath-q_\jmath|\leq h_\jmath/4\},
 \qquad
 v_x:=\sqrt{p_x\tau}+\tau.                                       \tag{6}
\]
Since \(h_\jmath=H_0\{\sqrt{c_\jmath\tau}+\tau\}>0\), every radius is positive.  Directly from the truncation at zero,
\[
 h_\jmath/2\leq r_\jmath\leq h_\jmath,
 \qquad |b_\jmath-c_\jmath|\leq h_\jmath/2.                     \tag{7}
\]
On \(G_x\), if \(c_\jmath\geq h_\jmath\), then \(b_\jmath=c_\jmath\) and \(|q_\jmath-b_\jmath|\leq r_\jmath\).  If \(c_\jmath<h_\jmath\), then \(\ell_\jmath=0\), while nonnegativity and \(q_\jmath\leq c_\jmath+h_\jmath/4<u_\jmath\) again give \(q_\jmath\in[\ell_\jmath,u_\jmath]\).  Hence
\[
 q_x\in Q_x\quad\hbox{on }G_x.                                  \tag{8}
\]

The inequality \(c\leq q+(H_0/4)(\sqrt{c\tau}+\tau)\) on \(G_x\), viewed as a quadratic inequality in \(\sqrt c\), implies \(c\leq C(q+\tau)\).  Substitution back into the reverse inequality gives \(q\leq C(c+\tau)\).  Therefore (7), \(\sum_\jmath q_\jmath=p_x\), and Cauchy--Schwarz in four coordinates imply
\[
 \sum_\jmath r_\jmath
 \leq C\left\{\sqrt{\tau}\sum_\jmath\sqrt{q_\jmath}+\tau\right\}
 \leq C v_x.                                                       \tag{9}
\]
Moreover, \(r_\jmath^2\geq(H_0^2/4)(c_\jmath\tau+\tau^2)\), so
\[
 \frac{q_\jmath}{m r_\jmath^2}
 \leq \frac{C(c_\jmath+\tau)}{m\tau(c_\jmath+\tau)}
 =\frac C L.                                                       \tag{10}
\]
No division by a cell probability was used.  In particular, if \(q_\jmath=0\), then its pilot and evaluation counts are zero almost surely, whereas its radius remains positive because \(\tau>0\).

We now control \(G_x^c\) at the same local scale.  If \(W\sim\operatorname{Pois}(\lambda)\), exponential Markov applied to
\[
 E e^{t(W-\lambda)}=\exp\{\lambda(e^t-1-t)\}
\]
and the corresponding expression for \(-t\) gives, for \(z\geq1\),
\[
 \Pr\{|W-\lambda|>\sqrt{2\lambda z}+2z\}\leq2e^{-z}.             \tag{11}
\]
Indeed, optimizing the Chernoff exponent and using \(e^t-1-t\leq t^2/[2(1-t/3)]\) for \(0<t<3\) gives the upper tail, and \(e^{-t}-1+t\leq t^2/2\) gives the lower tail.  Choose the fixed constant \(H_0\) sufficiently large.  Splitting into \(\lambda\geq C L\) and \(\lambda<C L\) shows that
\[
 |W-\lambda|\leq\sqrt{200\lambda L}+200L
 \quad\Longrightarrow\quad
 |W/m-\lambda/m|\leq\frac{H_0}{4}
 \{\sqrt{(W/m)\tau}+\tau\}.                                    \tag{12}
\]
For the first case, (12)'s premise gives \(W\geq\lambda/2\) after enlarging \(C\), so the square-root term controls \(\sqrt{\lambda L}/m\); for the second case the additive \(L/m\) term controls both \(\lambda/m\) and \(W/m\).  Thus the coordinate pilot-bad event is contained in the tail in (11) with \(z=100L\).

For completeness, integrating (11) over \(z\geq100L\), after substituting the threshold \(\sqrt{2\lambda z}+2z\), gives for \(s\in\{1,2,4\}\)
\[
 E\!\left[
 \left\{|W/m-\lambda/m|+\sqrt{(W/m)\tau}+\tau\right\}^{s}
 \mathbf 1_{\mathrm{bad}}
 \right]
 \leq C e^{-40L}\left\{\sqrt{(\lambda/m)\tau}+\tau\right\}^{s}. \tag{13}
\]
To verify the scale in (13), use \(\sqrt W\leq\sqrt\lambda+\sqrt{|W-\lambda|}\) and the layer-cake identity
\[
 E[R^s\mathbf 1\{R>r_0\}]
 =r_0^s\Pr(R>r_0)+\int_{r_0}^{\infty}s r^{s-1}\Pr(R>r)\,dr.
\]
For \(z\geq100L\), each fixed power of the ratio of \(\sqrt{\lambda z}+z\) to \(\sqrt{\lambda L}+L\) is at most a fixed polynomial in \(z/L\), whose integral against \(e^{-z}\) is at most \(C e^{-40L}\).  The identical calculation beginning at \(z=1\), together with the bounded central region, gives the unconditional version of (13) without the indicator and without the exponential factor.

The deterministic inequalities in (7) show that
\[
 \|b_x-q_x\|_1+\sum_\jmath r_\jmath
 \leq C\sum_\jmath
 \{|c_\jmath-q_\jmath|+\sqrt{c_\jmath\tau}+\tau\}.              \tag{14}
\]
The four pilot coordinates are independent by (1).  Apply a union bound to \(G_x^c\), use (13) on a bad coordinate, and use the unconditional moment bound following (13) on the other coordinates.  Cauchy--Schwarz with the fourth moments when \(s=2\) loses at most half the exponential exponent.  Since \(\sum_\jmath\sqrt{q_\jmath}\leq2\sqrt{p_x}\), this proves
\[
 E\!\left[
 \left\{\|b_x-q_x\|_1+\sum_\jmath r_\jmath\right\}^{s}
 \mathbf1_{G_x^c}\right]
 \leq C e^{-20L}v_x^s,
 \qquad s=1,2.                                                      \tag{15}
\]
Equation (9) also gives
\[
 E\!\left[\left(\sum_\jmath r_\jmath\right)^2
 \mathbf1_{G_x}\right]\leq C v_x^2.                              \tag{16}
\]
Thus exceptional pilots retain the local \(p_x\)-dependent scale.

We turn to approximation.  On \(G_x\), apply \(\mathrm{lem:simultaneous\mbox{-}jackson\mbox{-}certificate}\) to \(q_x\in Q_x\).  If \(\ell_\jmath=0\), then \(c_\jmath\leq h_\jmath\), and solving this inequality gives \(c_\jmath+h_\jmath\leq C\tau\).  Hence
\[
 \sqrt{(q_\jmath-\ell_\jmath)(u_\jmath-q_\jmath)}
 \leq C\sqrt{q_\jmath\tau},
 \qquad r_\jmath\leq C\tau.                                    \tag{17}
\]
If \(\ell_\jmath>0\), then \(c_\jmath>h_\jmath\); on \(G_x\), \(q_\jmath\geq3c_\jmath/4\) and \(r_\jmath=h_\jmath\leq C\sqrt{q_\jmath\tau}\).  Combining the two cases, the pointwise Jackson inequality and \(\sum_\jmath\sqrt{q_\jmath}\leq2\sqrt{p_x}\) give
\[
 |P_{x,K_d}(q_x)-\bar f_\epsilon(q_x)|
 \leq C_\epsilon\left\{
 \frac{\sqrt{p_x\tau}}{K_d}+\frac{\tau}{K_d^2}\right\}.         \tag{18}
\]
The same polynomial supplies the coefficient envelope required below; no second approximant is introduced.

Conditional on the pilot and on \(G_x\), apply (5) with \(t=h\), \(z=b_\jmath\), (8), and (10).  For every nonnegative integer \(h\),
\[
 \begin{split}
 \frac{E\{U_h(N_{\jmath,x};b_\jmath)^2\}}{r_\jmath^{2h}}
 &\leq\sum_{\ell=0}^h\binom h\ell^2\ell!(C/L)^\ell\\
 &\leq\sum_{\ell=0}^{\infty}\frac{(Ch^2/L)^\ell}{\ell!}
 =\exp(Ch^2/L),
 \end{split}                                                       \tag{19}
\]
where \(\binom h\ell\leq h^\ell/\ell!\).  Factorial unbiasedness also gives
\[
 E(Z_x\mid N')=P_{x,K_d}(q_x).                                    \tag{20}
\]
Write
\[
 P_{x,K_d}(v)-\bar f_\epsilon(b_x)
 =\sum_\alpha b_{x,\alpha}\{(v-b_x)/r_x\}^\alpha,
 \qquad
 S_x=(1+\epsilon^{-1})\sum_\jmath r_\jmath.
\]
On \(G_x\), evaluation counts are independent across coordinates by (1).  Minkowski's inequality in \(L_2\), (19), and the coefficient envelope in \(\mathrm{lem:simultaneous\mbox{-}jackson\mbox{-}certificate}\) therefore imply
\[
 E\!\left[\{Z_x-\bar f_\epsilon(b_x)\}^2\mid N'\right]
 \leq S_x^2\exp(C_1K_d+C_2K_d^2/L)
 \quad\hbox{on }G_x.                                               \tag{21}
\]
Indeed each coordinate degree is at most \(2(K_d-1)\), so \(\sum_\jmath\alpha_\jmath^2\leq16K_d^2\); (19) controls the \(L_2\)-norm of each tensor factorial monomial, and summing those norms costs exactly the declared coefficient \(\ell_1\)-envelope.

Choose the fixed \(\kappa>0\) sufficiently small and then the fixed cutoff \(D_0\) sufficiently large, as permitted by their declarations.  Since \(K_d=\max\{2,\lfloor\kappa L\rfloor\}\), throughout the Jackson branch
\[
 \exp(C_1K_d+C_2K_d^2/L)\leq C d^{1/16}.                           \tag{22}
\]
Increasing \(C\) handles the finitely many values absorbed into the cutoff.

It remains to verify clipping.  Put \(W_x=Z_x-\bar f_\epsilon(b_x)\) and \(t_x=d^{1/4}S_x>0\).  For every \(w\in\mathbb R\) and \(t>0\), direct consideration of \(|w|\leq t\) and \(|w|>t\) gives
\[
 |\Pi_{[-t,t]}(w)-w|\leq w^2/t,
 \qquad |\Pi_{[-t,t]}(w)|\leq|w|.                                \tag{23}
\]
On \(G_x\), equations (18), (20)--(23) imply
\[
 \begin{split}
 |E(T_x^{\mathrm{JF}}\mid N')-\bar f_\epsilon(q_x)|
 &\leq C_\epsilon\left\{
 \frac{\sqrt{p_x\tau}}{K_d}+\frac{\tau}{K_d^2}\right\}
   +CS_xd^{-3/16}.
 \end{split}                                                       \tag{24}
\]
For \(d\geq D_0\), \(K_d\geq cL\) and \(d^{-3/16}\leq C/L^2\).  Using (9) and \(v_x/L=\sqrt{p_x/(mL)}+1/(mL)\), (24) becomes
\[
 |E(T_x^{\mathrm{JF}}\mid N')-\bar f_\epsilon(q_x)|
 \leq C_\epsilon\left\{\sqrt{\frac{p_x}{mL}}+\frac1{mL}\right\}
 \quad\hbox{on }G_x.                                              \tag{25}
\]

The global Lipschitz conclusion of \(\mathrm{prop:identification\mbox{-}and\mbox{-}extension}\), (8), and the definition of \(S_x\) give \(|\bar f_\epsilon(b_x)-\bar f_\epsilon(q_x)|\leq S_x\) on \(G_x\).  Hence (21)--(23) also give
\[
 E\!\left[\{T_x^{\mathrm{JF}}-\bar f_\epsilon(q_x)\}^2
 \mid N'\right]\mathbf1_{G_x}
 \leq C S_x^2d^{1/16}\mathbf1_{G_x}.                              \tag{26}
\]
On \(G_x^c\), clipping and the same global Lipschitz bound yield the deterministic inequality
\[
 |T_x^{\mathrm{JF}}-\bar f_\epsilon(q_x)|
 \leq(1+\epsilon^{-1})\|b_x-q_x\|_1+d^{1/4}S_x.                  \tag{27}
\]
By (15), the first moment of the right side on \(G_x^c\) is at most \(C_\epsilon d^{1/4}e^{-20L}v_x\), and its second moment is at most \(C_\epsilon d^{1/2}e^{-20L}v_x^2\).  Since
\[
 d^{1/4}e^{-20L}\leq C/L,
 \qquad d^{1/2}e^{-20L}\leq C d^{1/16},                            \tag{28}
\]
these are precisely bounded by the right side of (25) and by \(C_\epsilon d^{1/16}v_x^2\), respectively.  Thus (15) explicitly proves that pilot failure has the asserted local scale.

Finally integrate (25)--(27) over the pilot.  Equation (16) and
\[
 v_x^2=(\sqrt{p_xL/m}+L/m)^2
 \leq2\{p_xL/m+L^2/m^2\}                                        \tag{29}
\]
give
\[
 \left|E_{\mathbb P}T_x^{\mathrm{JF}}-\bar f_\epsilon(q_x)\right|
 \leq C_\epsilon\left\{\sqrt{\frac{p_x}{mL_d}}+\frac1{mL_d}\right\} \tag{30}
\]
and
\[
 \operatorname{Var}_{\mathbb P}(T_x^{\mathrm{JF}})
 \leq E_{\mathbb P}\{T_x^{\mathrm{JF}}-\bar f_\epsilon(q_x)\}^2
 \leq C_\epsilon d^{1/16}
 \left\{\frac{p_xL_d}{m}+\frac{L_d^2}{m^2}\right\}.             \tag{31}
\]
All constants in (30)--(31) are uniform in \(\mathbb P\in\mathcal D_{d,\epsilon}^{\mathrm{obs}}\) and in the cell \(x\); their only nonuniversal dependence is on the fixed overlap constant \(\epsilon\).  The argument used neither a causal completion nor positivity of \(p_x\), an outcome mean, or a treatment-effect contrast, so it includes null cells, outcome-boundary means, unequal propensities, and exact ties.